
\[
\begin{tikzcd}[column sep=large,row sep=large]
V[T_1,\ldots,T_n] \arrow[r,"j'"] \arrow[d] & Z[T'_1,\ldots,T'_m,T_1,\ldots,T_n] \arrow[d]\\
V \arrow[r,"j"] & Z[T'_1,\ldots,T'_m],
\end{tikzcd}
\]
où $j'$ est obtenu à partir de $j$ par changement de base ; ce changement de base étant plat, $j'$ est une immersion régulière, et par suite, $j'\circ i$ est une immersion régulière de $U$ dans un schéma lisse sur $W$, d'où le résultat.

\textbf{Proposition 1.6.} \uline{Soit} $f:X\to Y$ \uline{un morphisme} $g:Y'\to Y$ \uline{un changement de base plat}. \uline{Si} $f$ \uline{est d'intersection complète}, \uline{alors le morphisme} $f':X\times_Y Y'\to Y'$ \uline{déduit de} $f$ \uline{l'est aussi}. \uline{La réciproque est vraie si} $g$ \uline{est de plus quasi-compact et surjectif}.

On peut trouver un recouvrement de $X$ par des ouverts $U_i$ tels que pour tout $i$, la restriction de $f$ à $U_i$ se factorise en $U_i\longrightarrow V_i\longrightarrow Y$, où $U_i\to V_i$ est une immersion régulière, et $V_i\to Y$ un morphisme lisse. Les $U_i\times_Y Y'$ forment un recouvrement ouvert de $X\times_Y Y'$, les immersions $U_i\times_Y Y'\longrightarrow V_i\times_Y Y'$ sont régulières d'après VII 1.5, et les morphismes $V_i\times_Y Y'\to Y'$ sont lisses, d'où le premier résultat. Pour la réciproque on conclut d'abord de EGA IV 2.7.1 (iv) que $f$ est localement de présentation finie, ce qui nous ramène facilement, grâce à 1.2, au cas ou $f$ est une immersion fermée, auquel cas l'assertion n'est autre que VII 1.5.

\textbf{Proposition 1.7.} \uline{Soit} $f:X\to Y$ \uline{un morphisme d'intersection complète}. \uline{Alors} $f$ \uline{est un morphisme parfait}.

La propriété étant locale sur $X$, on peut se placer au voisinage d'un point $x$ de $X$, et supposer qu'il existe un ouvert $U$ contenant $x$, et une factorisation :
